\documentclass[10pt,a4paper]{amsart}
\usepackage[margin=19mm]{geometry}
\usepackage{amsmath,amssymb,mathtools}
\usepackage[scr=rsfs,cal=euler]{mathalfa}
\usepackage{xfrac}
\usepackage[unicode,hidelinks,bookmarks=false]{hyperref}
\newcommand{\ie}{i.e.\ }
\newcommand{\cf}{cf.\ }
\newcommand{\resp}{respectively\ }
\newcommand{\Subsection}{\subsection}
\providecommand{\nobreakdash}{\nobreak\hbox{-}}
\setlength{\emergencystretch}{2em}
\multlinegap0pt
\usepackage{bm}
\usepackage{bbm}
\usepackage{dsfont}
\usepackage{xspace}
\usepackage{cancel}
\usepackage{mathtools}
\usepackage{amsthm}
\makeatletter
\@ifundefined{definition}{\newtheorem{definition}{Definition}}{}
\@ifundefined{lemma}{\newtheorem{lemma}[definition]{Lemma}}{}
\makeatother
\title[$q$-analogs of rational numbers: from Ostrowski numeration s]{$q$-analogs of rational numbers: from Ostrowski numeration systems to perfect matchings}
\author{Jean-Christophe Aval, S\'ebastien Labb\'e}
\date{Source: arXiv:2511.11290v1 (CC BY 4.0)}
\begin{document}
\maketitle

\noindent\emph{Source and licence.} $q$-analogs of rational numbers: from Ostrowski numeration systems to perfect matchings. Version arXiv:2511.11290v1 (CC BY 4.0), distributed under the Creative Commons Attribution 4.0 International licence, as recorded in the arXiv licence metadata for this exact version and shown on the abstract page https://arxiv.org/abs/2511.11290v1. Full rights notice: licenses/arxiv-2511.11290-CC-BY-4.0.txt.\par\smallskip

\noindent\textit{Editorial scope.} Counted unit: the Lemma whose source label is \texttt{lem:existence-rep-in-interval} in the article (source0.tex lines 1419--1438, source label \texttt{lem:existence-rep-in-interval}), delivered together with the article's own complete original proof of it (source0.tex lines 1440--1602, one \texttt{proof} environment of 163 lines). No mathematics is written, repaired or re-derived: statement and proof are the article's own text line for line. Required context retained uncounted: lem:upper-bounds-on-the-sum (lines 1388--1401). Every definition, display and label of those lines is retained unchanged. Omitted from this package and not redistributed: 1--1387, 1402--1418, 1439--1439, 1603--5699. Nothing inside a retained excerpt is omitted or altered: every retained body is delivered exactly as the article's lines are written, so no omission receipt of any kind accompanies this package. Citations: the retained excerpts carry no citation command at all, so no bibliography entry of the article is transcribed and no reference is redistributed. Reference adaptation: none was needed and none was made; every reference inside the delivered unit resolves inside it, so no unresolved reference or citation is produced. Printed slips kept literal: no printed word, symbol, hypothesis or number of the counted statement or of its proof was altered. Layout and class: the article's own class is \texttt{amsart} with packages including amsmath, amsfonts, amssymb, amsthm, geometry, latexsym, graphicx, url, enumerate, cite, todonotes, hyperref, tikz, tikz-cd; the wrapper is the lane's pinned \texttt{amsart} editorial wrapper at 10pt on A4, which restates the theorem-like environments the retained text needs and adds the article's own macro definitions that the retained text needs (none), with extra wrapper packages (bm, bbm, dsfont, xspace, cancel, mathtools). The delivered numbering is the unit's own. Assets: no retained span contains a figure environment, a graphics-inclusion command, a PDF-page inclusion, a table or a TikZ picture; no image file is copied into this package, the package declares no asset dependency and no image file is redistributed with it. Licence: version arXiv:2511.11290v1 of this article is distributed under the Creative Commons Attribution 4.0 International licence, as recorded in the arXiv licence metadata for this exact version and shown on the abstract page https://arxiv.org/abs/2511.11290v1. A full rights notice is delivered at licenses/arxiv-2511.11290-CC-BY-4.0.txt. Characters: every retained excerpt is pure ASCII, so the delivered engine, XeLaTeX with its default Latin Modern text font, reports no missing character.
\begin{lemma}\label{lem:upper-bounds-on-the-sum}
    Let $a=[a_0; a_1, \dots, a_{k-1}]$ be the continued fraction expansion
    of a positive rational number.
    Let $n= \sum_{i=0}^{k-1} (-1)^{i}b_{i}r_{i}$
    where $(b_i)_{0\leq i< k}$ is an admissible
    sequence for $a$. Then,
    \[
        n <
        \begin{cases}
            r_{k-1} & \text{ if $k$ is even},\\
            r_k     & \text{ if $k$ is odd}.
        \end{cases}
    \]
\end{lemma}

\begin{lemma}\label{lem:existence-rep-in-interval}
    Let $[a_0; a_1, \dots, a_{k-1}]$ be the continued fraction expansion
    of a positive rational number.
    \begin{enumerate}
        \item If $k$ is odd, then every integer $n$ in the interval
                $0\leq n< r_{k}$
                can be written as
                $
                    n = \sum_{i=0}^{k-1} (-1)^{i}b_{i}r_{i}
                $
                where $(b_i)_{0\leq i\leq k-1}$ is an admissible sequence.
        \item If $k$ is even, then every integer $n$ in the interval
                $r_{k-1}-r_{k}\leq n< r_{k-1}$
                can be written as
                $
                    n = \sum_{i=0}^{k-1} (-1)^{i}b_{i}r_{i}
                $
                where $(b_i)_{0\leq i\leq k-1}$ is an admissible sequence.
    \end{enumerate}
\end{lemma}

\begin{proof}
    The proof of existence of the admissible sequence is done by induction.
    We prove that if (1) holds for some odd $k$, then (2) hold for $k+1$,
    and that if (2) holds for some even $k$, then (1) hold for $k+1$.

    First, we prove that (1) holds for $k=1$.
    Recall that we have $r_{-1}=r_0=1$.
    Let $k=1$ and $n$ be an integer in the interval $0\leq n<r_{k}=r_1=a_0r_0+r_{-1}=a_0+1$.
    Let $b_0=n$. We have $0\leq b_0\leq a_0$. Thus, the sequence $(b_0)$ is admissible
    and satisfies
    \[
        n = b_0 = (-1)^{0}b_0r_0.
    \]

    Suppose that (1) holds for some odd integer $k\geq1$.
    We want to show that (2) holds for $k+1$.
    Let $n$ be an integer in the interval
    $r_{k}-r_{k+1}\leq n< r_{k}$.
    Let $Q$ be the quotient and $R$ be the remainder
    of the division of $n$ by $r_{k}$.
    We have
    \begin{align*}
        0
        &=    \left\lfloor
        \frac{r_{k}-1}
             {r_{k}}
            \right\rfloor
        \geq Q
        = \left\lfloor
        \frac{n}
             {r_{k}}
            \right\rfloor
        \geq
            \left\lfloor
        \frac{r_{k}-r_{k+1}}
             {r_{k}}
            \right\rfloor\\
        &=
            \left\lfloor
        \frac{r_{k}-(a_{k}r_{k}+r_{k-1})}
             {r_{k}}
            \right\rfloor
        = - a_{k}
            +
            \left\lfloor
        \frac{r_{k}-r_{k-1}}
             {r_{k}}
            \right\rfloor
        = - a_{k}.
    \end{align*}
    Let $b_{k}=-Q$ so that
    $0\leq b_{k}\leq a_{k}$.
    By definition of the division, we have
    $n=Qr_{k} + R$
    where the remainder satisfies $0\leq R<r_{k}$.
    From (1), the integer $R$
    can be written uniquely as
    $
        %
        \sum_{i=0}^{k-1} (-1)^{i}b_{i}r_{i}
    $
    where $(b_i)_{0\leq i\leq k-1}$ is an admissible sequence.
    Thus,
    \begin{align*}
        n
        &=Qr_{k} + R
         =-b_{k}r_{k} + \sum_{i=0}^{k-1} (-1)^{i}b_{i}r_{i}
         =\sum_{i=0}^{k} (-1)^{i}b_{i}r_{i}.
    \end{align*}
    It remains to show that $(b_i)_{0\leq i\leq k}$ is an admissible sequence,
    that is, that $b_{k}=a_{k}$ implies that $b_{k-1}=a_{k-1}$.
    Suppose that $b_{k}=a_{k}$. Thus, $Q=-a_{k}$.
    We have
    \[
        r_{k}
        > R
        =
        n - Q r_{k}
        =
        n+a_{k}r_{k}
        \geq
        r_{k}-r_{k+1} +a_{k}r_{k}
        =
        r_{k}-r_{k-1}.
    \]
    Let $R'=\sum_{i=0}^{k-2} (-1)^{i}b_{i}r_{i}$.
    From Lemma~\ref{lem:upper-bounds-on-the-sum}, we have $R'<r_{k-2}$ since $k-1$ is even.
    Therefore,
    \[
        b_{k-1}r_{k-1}
        = R - R'
        > r_{k}-r_{k-1} - r_{k-2}
        = a_{k-1}r_{k-1}+r_{k-2}-r_{k-1} - r_{k-2}
        = (a_{k-1}-1)r_{k-1}.
    \]
    This implies that $b_{k-1}=a_{k-1}$ in the sum $R=\sum_{i=0}^{k-1} (-1)^{i}b_{i}r_{i}$.
    We conclude that (2) holds for $k+1$.

    Suppose that (2) holds for some even integer $k\geq1$.
    We want to show that (1) holds for $k+1$.
    Let $n$ be an integer in the interval
    $0\leq n< r_{k+1}$.
    Let $b_{k}$ be the quotient and $R$ be the remainder
    of the division of $n-(r_{k-1}-r_{k})$ by $r_{k}$.
    We have
    \begin{align*}
        0\leq b_{k} &= \left\lfloor
        \frac{n-(r_{k-1}-r_{k})}
             {r_{k}}
            \right\rfloor
        \leq
            \left\lfloor
        \frac{r_{k+1}{-}1-(r_{k-1}-r_{k})}
             {r_{k}}
            \right\rfloor\\
        &=
            \left\lfloor
        \frac{a_{k}r_{k}+r_{k-1}-(r_{k-1}-r_{k}){-}1}
             {r_{k}}
            \right\rfloor
        =
            \left\lfloor
        \frac{a_{k}r_{k}+r_{k}-1}
             {r_{k}}
            \right\rfloor
        = a_{k}.
    \end{align*}
    By definition of the division, we have
    $n-(r_{k-1}-r_{k})=b_{k}r_{k} + R$
    where the remainder satisfies $0\leq R<r_{k}$.
    Thus,
    \[
        r_{k-1}-r_{k} \leq R+ r_{k-1}-r_{k}  < r_{k-1}.
    \]
    From (2), the integer $R+ r_{k-1}-r_{k}$
    can be written uniquely as
    $
        %
        \sum_{i=0}^{k-1} (-1)^{i}b_{i}r_{i}
    $
    where $(b_i)_{0\leq i\leq k-1}$ is an admissible sequence.
    Therefore,
    \[
        n
        = b_{k}r_{k} + R+(r_{k-1}-r_{k})
        = b_{k}r_{k} + \sum_{i=0}^{k-1} (-1)^{i}b_{i}r_{i}
        = \sum_{i=0}^{k} (-1)^{i}b_{i}r_{i}.
    \]
    It remains to show that $(b_i)_{0\leq i\leq k}$ is an admissible sequence,
    that is, that $b_{k}=0$ implies that $b_{k-1}=0$.
    Suppose that $b_{k}=0$.
    We have
    $n = -b_{k-1}r_{k-1} + R'$
    where $R'=\sum_{i=0}^{k-2} (-1)^{i}b_{i}r_{i}$.
    From Lemma~\ref{lem:upper-bounds-on-the-sum}, we have $R'<r_{k-1}$ since $k-1$ is odd.
    Therefore, using $n\geq0$,
    \[
        b_{k-1}r_{k-1} = R' - n < r_{k-1} + 0 = r_{k-1}.
    \]
    Thus, we conclude that $b_{k-1}=0$
    and $(b_i)_{0\leq i\leq k}$ is an admissible sequence.
    We conclude that (1) holds for $k+1$.
\end{proof}

\end{document}
